problem:  
Let \((M^n,g)\) be a complete noncompact Riemannian manifold satisfying a volume doubling condition and a scale-invariant Poincaré inequality, with Ricci curvature bounded below by  
\[
\mathrm{Ric}_g \ge -K(1+r(x)^2).
\]  
For \(k>0\), denote by \( \mathcal{H}_k(M)\) the space of harmonic functions \(u\) satisfying the growth bound \(|u(x)| \le C(1+r(x))^k\).  By results of Colding–Minicozzi and Li–Wang, if \(\mathrm{Ric}\ge0\) then \(\dim \mathcal{H}_k(M) < \infty\) and grows at most polynomially in \(k\). Little is known, however, when Ricci curvature decays quadratically and volume doubling/Poincaré still hold.

**Open problem:**  
Does there exist a sharp quantitative relationship between the rate of Ricci curvature decay and the dimension growth of polynomially bounded harmonic functions?  
More precisely, for manifolds satisfying  
\[
\mathrm{Ric}_g(x) \ge -K(1+r(x)^2),
\]
is it true or false that there exists an exponent \(\alpha=\alpha(n,K,C_D,C_P)\) such that  
\[
\dim \mathcal{H}_k(M) \le C(1+k)^\alpha \quad \forall k>0,
\]
and that this bound is sharp—in the sense that one can construct examples (with the same curvature decay rate) achieving equality in the asymptotic order of \((1+k)^{\alpha}\)?  

Equivalently, does quadratic Ricci decay admit a universal *spectral-type polynomial bound* on the space of harmonic functions of polynomial growth under volume doubling and Poincaré assumptions, analogous to the Colding–Minicozzi finiteness results for nonnegative Ricci curvature?  

Why is it a "good" problem:  
This problem is solidly grounded in the structure of geometric analysis: it asks for the precise link between curvature decay, analytic inequality structure (doubling and Poincaré), and the size of the space of harmonic functions of controlled growth. It moves beyond the classical Liouville question into dimensional asymptotics—connecting geometric spectral properties, analytic capacity bounds, and asymptotic geometry. Establishing or disproving a universal growth law would generalize the deep finiteness theorems under nonnegative Ricci curvature to a far weaker geometric regime, pushing the understanding of vanishing and rigidity phenomena into the frontier where stochastic and PDE methods overlap in geometric analysis. The question is simple to state but could only be resolved through a synthesis of curvature analysis, harmonic measure estimates, and heat kernel asymptotics—making it a “narrow entrance, profound depth” research problem.